% ============================================================================
% Cover Letter — Paper 5 (JoC / BMC)
% Target: Journal of Cheminformatics
% ============================================================================
\documentclass[11pt,a4paper]{letter}

\usepackage[margin=2.3cm]{geometry}
\usepackage{microtype}
\usepackage{hyperref}
\usepackage{enumitem}
\usepackage{siunitx}
\sisetup{
  detect-all           = true,
  group-minimum-digits = 4,
}

\signature{Myke Vital Sao Temgoua \\ \small Corresponding Author}
\address{%
  University of Yaound\'{e} I \\
  Department of Physics, Faculty of Science \\
  P.O. Box 812, Yaound\'{e}, Cameroon \\
  \href{mailto:myke-vital.sao@facsciences-uy1.cm}{myke-vital.sao@facsciences-uy1.cm}
}

\begin{document}

\begin{letter}{%
  The Editors \\
  \textit{Journal of Cheminformatics} \\
  BioMed Central
}

\opening{Dear Editors,}

We are pleased to submit our manuscript, \textbf{``Larger models, unchanged verdict: a controlled benchmark of graph neural networks and transformers against ECFP4 fingerprints for antimalarial natural-product activity prediction''}, for consideration as a Research Article in the \textit{Journal of Cheminformatics}.

The ``scale-centric'' narrative in molecular machine learning holds that larger pretrained models should supersede compact cheminformatics baselines. We test this assumption on a frozen panel of \num{19836} African antimalarial natural products under both random and Bemis--Murcko scaffold cross-validation (\num{25} fold--seed replicates), comparing a random forest on ECFP4 fingerprints against a graph isomorphism network (GIN), two topological-fusion variants (GIN--TFP, GIN--TNE), and a fine-tuned ChemBERTa transformer.

Under the scaffold split, the practically relevant out-of-distribution setting, \textbf{ECFP4--RF achieves the highest ROC AUC (\num{0.8300})}, ahead of GIN--TFP (\num{0.8138}), GIN--TNE (\num{0.8090}), GIN (\num{0.8047}) and ChemBERTa (\num{0.7867}); no graph or transformer arm significantly outperforms the fingerprint baseline. Under the random split the verdict is identical: every arm is significantly below ECFP4--RF (\num{0.9433}) --- GIN--TFP \num{0.9084}, GIN--TNE \num{0.8918}, GIN \num{0.9098}, ChemBERTa \num{0.9121} (all paired $p < \num{0.0001}$, Benjamini--Hochberg adjusted) --- so the negative result holds under both in-distribution and out-of-distribution evaluation. We document and correct a subtle cross-fold weight-leakage artifact in transformer fine-tuning that would otherwise inflate reported AUCs by up to \num{0.15}, and an attribution analysis shows that persistent-image topological features carry the majority of the predictive signal, linking learned representations to molecular geometry even when predictive gains are absent.

We believe this controlled, leak-audited, fully reproducible benchmark aligns closely with the scope of the \textit{Journal of Cheminformatics}---including its collection on the evaluation of AI and machine learning models in cheminformatics---and contributes an honest-negative result on a clinically motivated natural-product panel, mirroring the emerging consensus that representation--task alignment, not model scale, governs small-molecule activity prediction. Code, frozen splits and trained checkpoints are released under an open licence and are available to reviewers in the public repository (\url{https://github.com/NanaEngo/Malaria_codesV2}).

This work is original, is not under consideration elsewhere, and all authors have approved the submission.

\closing{We thank you for your consideration.}

\ps{%
\textbf{Suggested reviewers:}\\
\begin{itemize}[leftmargin=*, itemsep=1pt]
  \item Prof.\ Gisbert Schneider (generative molecular design, fingerprints)
  \item Dr.\ Ola Engkvist (de novo design, graph models)
  \item Dr.\ Davide Boldini (natural-product fingerprints and chemical space)
\end{itemize}
}

\end{letter}
\end{document}
